From Thomson's formula $f = \frac{1}{2\pi\sqrt{L\,C}}$:
\begin{abcliste}
\abc
\begin{align}
 C_1 &= \CaF = \frac{1}{4\pi^2\times\left(\fa\right)^2\times\Li} = \Ca \approx \result{\CaP{-12}{3}}
\end{align}

\abc With $L$ unchanged, $f \propto 1/\sqrt{C}$, so $C \propto 1/f^2$:
\begin{align}
 \frac{C_2}{C_1} &= \kCF = \left(\frac{\fa}{\fb}\right)^2 = \kC \approx \result{\kCP{0}{3}} \\
 C_2 &= \CbF = \Cb \approx \result{\CbP{-12}{3}}
\end{align}
The capacitance must be made about $\kCP{0}{3}$ times as large.

\abc The classmate is wrong twice. In $L\,\ddot Q = -\frac{1}{C}\,Q$ the role of the spring constant is played by $1/C$: a $\result{\text{larger capacitance is a softer spring}}$, and a softer spring oscillates more slowly. To lower the frequency, $C$ must therefore be made larger, not smaller. And as $\omega = \sqrt{D/m}$, the frequency changes only with the square root of $1/C$: to lower it by the factor $\kwP{0}{3}$, $C$ must grow by the square of that factor, $\kCP{0}{3}$.
\end{abcliste}
